\section{Logarithmic speeds and continuation near the real axis}
\label{sec:common-regions}
The full logarithmic comparison gives counting information on the
sizes of the velocities. An absolutely convergent difference formula
then enlarges the continuation region in directions close to the real
parameter axis. Put $s_T=\sqrt{(\log\log T)/2}$.

\begin{theorem}\label{thm:gaussian-speed-component}
Fix a primitive character and a finite interval $I$ of positive length.
For either strict sign of $\Re v_\rho$, the lower proportion, relative
to $N_{\chi,T}$, of distinct simple centers for which any one of
\[
 \frac{\log|L_\chi'(\rho)|-\log\log T}{s_T},\qquad
 \frac{\log|v_\rho|-\log(T/\log T)}{s_T},\qquad
 \frac{\log|\Re v_\rho|-\log(T/\log T)}{s_T}
\]
belongs to $I$ is at least $10^{-16}\mathbb P(G\in I)$, where
$G$ is a standard real Gaussian. The constant is absolute; the
starting height may depend on the character and the interval.
\end{theorem}

\begin{proof}
Fix an auxiliary interval $J$ of Gaussian mass $p>0$. The five
interval limits give the same phase-discard margin as before.
Each strict sign within $X_\rho\in J$ has absolute $C_\rho$ mass
at least $pNL/2000000$, eventually.
Delete $|\cos\psi_\chi|<10^{-12}$. By the restricted count and
second moment its mass is at most
\[
 6\cdot10^{-12}
 \left(\#\{X_\rho\in J\}\sum_J|A_\rho|^2\right)^{1/2}
       \le (6\cdot10^{-12}\sqrt b+o(1))pNL.
\]
The unit endpoint strips cost $o(NL)$. Choose a fixed curvature
constant $K_J$ so large that
$6\sqrt{C_M}/K_J<p/16000000$. Formula
\eqref{eq:weighted-curvature-new} bounds the mass lost by
$M_\rho>K_JL|D_\rho|$ by this last quantity times $NL$.
The retained nonzero summands are simple and have separated centers,
as in Lemma~\ref{lem:large-derivative-selection}.

Apply the absolute logarithmic estimate
\eqref{eq:log-kernel-new}, its reversed mean
\eqref{eq:reversed-kernel-new}, and
\eqref{eq:prime-comparison-new} on these centers. For every fixed
$\epsilon>0$, only $o(N)$ retained centers have
\[
 |\log|L_\chi(\rho+\eta/L)|-\Re P_T(\rho)|
                                      >\epsilon\sigma_T.
\]
Indeed Chebyshev bounds the prime comparison exceptions and Markov
bounds the kernel exceptions. Choose $\eta$ fixed and small enough
for the curvature constant $K_J$. Taylor's formula gives
$L_\chi(\rho+\eta/L)=(\eta/L)D_\rho(1+e_\rho)$ with
$|e_\rho|\le1/8$. Its bounded logarithmic factors imply
\begin{equation}\label{eq:derivative-gaussian-transfer}
       Y_\rho:=\frac{\log|D_\rho|-\log L}{s_T}
                   =X_\rho+o(1)
\end{equation}
in counting measure on the retained set, because
$\sigma_T/s_T\to1$ and $J$ is bounded.
Any $o(N)$ exceptions have $o(NL)$ signed absolute mass, by
$|C_\rho|\le6|A_\rho|$ and the global second moment.
Each sign still has mass at least $pNL/4000000$.
The restricted squared mass is at most $72pNL^2$ eventually.
Cauchy's inequality therefore gives at least
\[
              \frac{p}{72\cdot4000000^2}N>10^{-16}pN
\]
centers of either sign after every deletion.
For the derivative assertion take an inner interval $J$ of $I$,
apply \eqref{eq:derivative-gaussian-transfer}, take the lower limit
in $T$ and then exhaust $I$. Gaussian endpoints have zero mass.

Finally the exact velocity formula and the exterior Euler bounds
give, uniformly on the dyadic interval,
\[
             \log|v_\rho|=\log T-\log|D_\rho|+O(1).
\]
The retained phase margin makes
$\log|\Re v_\rho|-\log|v_\rho|$ bounded as well.
The normalized logarithmic speeds are thus $-Y_\rho+o(1)$.
Apply the preceding argument to reflected inner intervals and use
the symmetry of the Gaussian distribution. The strict velocity sign
is the sign of $C_\rho$ at every retained center, proving all claims.
\end{proof}

\begin{theorem}\label{thm:anisotropic-continuation}
There is an absolute $c>0$ such that the same two sets of at least
$10^{-24}N_{\chi,T}$ centers in
Theorem~\ref{thm:common-deformation} have simple disjoint branches on
\[
 \Omega_T=\{u\in\mathbb C:
                 (1+T|u|)e^{12T|\Im u|}\le c e^{s_T}\}.
\]
The sections are holomorphic on the interior and continue locally
across every boundary point. Their ordinates remain in $(T,2T)$.
In particular the real interval $|u|\le c e^{s_T}/(2T)$ is included
eventually. The strict-side conclusion holds for $0<u\le c_0/T$.

More generally, for each fixed $A>0$ there are positive constants
$c_A,c'_A$ such that each strict initial sign has at least
\[
        10^{-17}\mathbb P\{G\in[A+1,A+2]\}N_{\chi,T}
\]
simple centers with these continuation and ordinate properties on
\[
                  (1+T|u|)e^{12T|\Im u|}\le c_A e^{2As_T},
\]
and with strict sides on $0<u\le c'_A/T$.
All constants are independent of the fixed character.
\end{theorem}

\begin{proof}
For $\Re s>0$ and $|u|<1/2$, the identity
\begin{equation}\label{eq:shift-difference-series}
 F_\chi(u,s)-L_\chi(s)
 =\sum_{n\ge1}\chi(n)n^{-s}\{(1+u/n)^{-s}-1\}
\end{equation}
is locally normally convergent. Each summand is
$O(|u|n^{-1-\Re s})$ on compact subsets. It equals the difference
of the original Dirichlet series in $\Re s>1$, and analytic
continuation proves the displayed identity. The logarithm is the
one on $\Re(1+u/n)>0$.

On a center disk $|s-\rho|\le h/L$, $h\le1/10$, put
$U=T|u|$, $v=T|\Im u|$ and $\sigma_s=\Re s$.
For $0\le t\le1$,
\[
 \log|1+tu/n|\ge-2|\Re u|/n,\qquad
 |\arg(1+tu/n)|\le2|\Im u|/n.
\]
Here $2/5\le\sigma_s\le3/5$ and $|s|,|\Im s|\le3T$ eventually.
The trivial difference estimate and the integral identity
\[
 (1+u/n)^{-s}-1=-\frac{su}n
                         \int_0^1(1+tu/n)^{-s-1}dt
\]
give
$|(1+u/n)^{-s}-1|\le3e^2e^{6v}\min(1,U/n)$.
Splitting the sum at $n=U$ yields $O(U)$ for $U\le1$ and
$O(U^{1-\sigma_s})$ for $U\ge1$.
Since $U<T$ and $|\sigma_s-1/2|\le h/L$, the last power differs
from $\sqrt U$ by a factor at most $e^h$.
Consequently the absolute constant $C_0=10^4$ suffices in
\begin{equation}\label{eq:anisotropic-perturbation}
 \sup_{|s-\rho|\le h/L}|F_\chi(u,s)-L_\chi(s)|
                \le C_0 e^{6T|\Im u|}\sqrt{1+T|u|}.
\end{equation}

For the original selected centers,
$\sigma_T-s_T=O_\chi((\log V)^2/\sqrt V)=o(1)$.
Their derivative floor therefore gives
$|D_\rho|h/L\ge d_0e^{s_T/2}$ with an absolute $d_0>0$,
eventually. Choose $c\le(d_0/(10C_0))^2$.
On $\Omega_T$, the perturbation in
\eqref{eq:anisotropic-perturbation} is at most
$|D_\rho|h/(10L)$, whereas the original Taylor remainder is at most
$K_0h|D_\rho|h/(2L)$. Rouch\'e with reference function
$D_\rho(s-\rho)$ gives exactly one zero counted with multiplicity
in each disk. It is simple. The strict boundary inequality persists
in a neighborhood of every parameter point, and uniqueness in each
fixed center disk patches the local sections. The regions are
eventually contained in $|u|<1/2$, so the proof applies throughout.
The disks already provide the distinctness and ordinate conditions.
The stated real interval lies in $\Omega_T$ for large $T$.

For the second assertion use the proof of
Theorem~\ref{thm:gaussian-speed-component} on $J=[A+1,A+2]$.
Remove the $o(N)$ centers where the transfer
\eqref{eq:derivative-gaussian-transfer} has error at least $1/2$.
The remaining signed mass and restricted squared mass still give
more than $10^{-16}\mathbb P(G\in J)N$ centers of each sign.
In particular the asserted smaller fraction has
\[
              |D_\rho|\ge Le^{As_T},\qquad
              M_\rho\le K_AL|D_\rho|.
\]
Choose their separated radius $h_A/L$ with $K_Ah_A\le1/100$.
With $c_A\le(h_A/(10C_0))^2$, the same Rouch\'e argument proves
the second region. The fixed phase margin and the local relative
velocity estimate in the preceding section give $c'_A/T$ after
choosing $c'_A$ small enough. No additional centers are discarded.
\end{proof}

For every fixed $B>0$, these regions include both the complex disk
$|u|\le B\sqrt{\log\log T}/T$ and the real interval
$|u|\le T^{-1}e^{B\sqrt{\log\log T}}$ for some positive fraction
of each initial sign. Indeed choose $A=12B+1$; the logarithms of the
two left sides are at most $12B\sqrt{\log\log T}+O_B(\log\log\log T)$
and $B\sqrt{\log\log T}+O(1)$, respectively, while
$2As_T=\sqrt2 A\sqrt{\log\log T}$.
The construction exhibits the same analytic zero sections on the
larger regions, with their real normal signs controlled on the
explicit interval supplied by the theorem.
